\begingroup
\begin{figure*}[ht]
\centering
\hspace{0.3in}
	\begin{subfigure}[b]{0.3\textwidth}
		\includegraphics[width=\textwidth]{CellSpansA.png}
		\subcaption[]{Corner Duplicates}
		\label{fig:CellSpansA}
	\end{subfigure}
\hspace{0.3in}
	\begin{subfigure}[b]{0.3\textwidth}
		\includegraphics[width=\textwidth]{CellSpansB.png}
		\subcaption[]{Collision Duplicates}
		\label{fig:CellSpansB}
	\end{subfigure}
\hspace{0.3in}
\caption[TITLE:Cell Spans]{Two-dimensional illustration of occupancy and collision duplication
arising from corner-based occupancy construction. (a) Multiple corners
may map to the same cell, resulting in duplicate occupancy
contributions. (b) A single collision may be represented within
multiple neighboring cells when interacting particles span a cell
boundary. The three-dimensional implementation follows identically
using the eight corners of the particle AABB.}
\Description[TITLE:Cell Spans]{Two-dimensional illustration of occupancy and collision duplication
arising from corner-based occupancy construction. (a) Multiple corners
may map to the same cell, resulting in duplicate occupancy
contributions. (b) A single collision may be represented within
multiple neighboring cells when interacting particles span a cell
boundary. The three-dimensional implementation follows identically
using the eight corners of the particle AABB.}
\label{fig:CellSpans}
\end{figure*}
\endgroup